\documentclass[12pt,aspectratio=169,ignorenonframetext]{ltxbeamer}

\usepackage[spanish,es-nodecimaldot,es-tabla]{babel}
\usepackage{tikz-feynman}
\usepackage{circuitikz}
\usepackage{asymptote}

\tikzfeynmanset{
    compat=1.1.0,
}

\ctikzset{
    open poles fill=black,
}

\sisetup{
    unit-mode=text,
    text-font-command=\lmr,
}

\title{Diagramación}
\subtitle{Curso de {\lmr\LaTeX}}
\author{Joar Buitrago\texorpdfstring{\thanks{jebuitragoc@unal.edu.co}}{ <jebuitragoc@unal.edu.co>}}
\date{}

\addbibresource{bibliography.bib}

\newcommand{\TikZ}{Ti\textit{k}Z}



\begin{document}
\begin{frame}
\maketitle
\end{frame}





\section{Diagramas de Feynman}

\begin{frame}
\frametitle{Diagramas de Feynman}

\begin{center}
    \begin{tikzpicture}
        \begin{feynman}[horizontal=a to b]
        	  \diagram {
                i1 -- [fermion] a -- [fermion] i2,
            	  a -- [photon] b,
            	  f1 -- [fermion] b -- [fermion] f2,
            };
        \end{feynman}
    \end{tikzpicture}
\end{center}
\end{frame}





\begin{frame}[fragile]
\frametitle{Diagramas de Feynman}

\begin{minted}{latex}
\usepackage{tikz}
\usepackage{tikz-feynman}
\tikzfeynmanset{
    compat=1.1.0,
}
\end{minted}
\end{frame}





\begin{frame}[fragile]
\frametitle{Diagramas de Feynman}

\begin{onlyenv}<1>
    \begin{minted}{latex}
\begin{tikzpicture}
    \begin{feynman}[horizontal=a to b]
        \diagram {
            i1 -- [fermion] a -- [fermion] i2,
            a -- [photon] b,
            f1 -- [fermion] b -- [fermion] f2,
        };
    \end{feynman}
\end{tikzpicture}
    \end{minted}
\end{onlyenv}

\begin{onlyenv}<2>
    \begin{center}
        \begin{tikzpicture}
            \begin{feynman}[horizontal=a to b]
                \diagram {
                    i1 -- [fermion] a -- [fermion] i2,
                    a -- [photon] b,
                    f1 -- [fermion] b -- [fermion] f2,
                };
            \end{feynman}
        \end{tikzpicture}
    \end{center}
\end{onlyenv}
\end{frame}





\begin{frame}[fragile]
\frametitle{Diagramas de Feynman}

\begin{onlyenv}<1>
    \begin{minted}[breaklines]{latex}
\begin{tikzpicture}
    \begin{feynman}[large, vertical=e to f]
        \diagram {
            a -- [fermion] b -- [photon, momentum=\(k\)] c -- [fermion] d,
            b -- [fermion, momentum'=\(p_{1}\)] e -- [fermion, momentum'=\(p_{2}\)] c,
            e -- [gluon]  f,
            h -- [fermion] f -- [fermion] i;
        };
    \end{feynman}
\end{tikzpicture}
    \end{minted}
\end{onlyenv}

\begin{onlyenv}<2>
    \begin{center}
        \begin{tikzpicture}
            \begin{feynman}[large, vertical=e to f]
                \diagram {
                	  a -- [fermion] b -- [photon, momentum=\(k\)] c -- [fermion] d,
                	  b -- [fermion, momentum'=\(p_{1}\)] e -- [fermion, momentum'=\(p_{2}\)] c,
                	  e -- [gluon]  f,
                	  h -- [fermion] f -- [fermion] i;
                };
            \end{feynman}
        \end{tikzpicture}
    \end{center}
\end{onlyenv}
\end{frame}





\begin{frame}[fragile]
\frametitle{Diagramas de Feynman}

\begin{onlyenv}<1>
    \tiny
    \begin{columns}
        \begin{column}{0.5\textwidth}
            \begin{minted}[breaklines]{latex}
\begin{tikzpicture}
    \begin{feynman}
        \vertex (a1) {\(\overline b\)};
        \vertex[right=1cm of a1] (a2);
        \vertex[right=1cm of a2] (a3);
        \vertex[right=1cm of a3] (a4) {\(b\)};
        \vertex[right=1cm of a4] (a5);
        \vertex[right=2cm of a5] (a6) {\(u\)};

        \vertex[below=2em of a1] (b1) {\(d\)};
        \vertex[right=1cm of b1] (b2);
        \vertex[right=1cm of b2] (b3);
        \vertex[right=1cm of b3] (b4) {\(\overline d\)};
        \vertex[below=2em of a6] (b5) {\(\overline d\)};

        \vertex[above=of a6] (c1) {\(\overline u\)};
        \vertex[above=2em of c1] (c3) {\(d\)};
        \vertex at ($(c1)!0.5!(c3) - (1cm, 0)$) (c2);
            \end{minted}
        \end{column}


        \begin{column}{0.5\textwidth}
            \begin{minted}[breaklines]{latex}
        \diagram* {
            {[edges=fermion]
                (b1) -- (b2) -- (a2) -- (a1),
                (b5) -- (b4) -- (b3) -- (a3) -- (a4) -- (a5) -- (a6),
            },
            (a2) -- [boson, edge label=\(W\)] (a3),
            (b2) -- [boson, edge label'=\(W\)] (b3),

            (c1) -- [fermion, out=180, in=-45] (c2) -- [fermion, out=45, in=180] (c3),
            (a5) -- [boson, bend left, edge label=\(W^{-}\)] (c2),
        };

        \draw [decoration={brace}, decorate] (b1.south west) -- (a1.north west) node [pos=0.5, left] {\(B^{0}\)};
        \draw [decoration={brace}, decorate] (c3.north east) -- (c1.south east) node [pos=0.5, right] {\(\pi^{-}\)};
        \draw [decoration={brace}, decorate] (a6.north east) -- (b5.south east) node [pos=0.5, right] {\(\pi^{+}\)};
    \end{feynman}
\end{tikzpicture}
            \end{minted}
        \end{column}
    \end{columns}
\end{onlyenv}

\begin{onlyenv}<2>
    \begin{center}
        \begin{tikzpicture}
            \begin{feynman}
                \vertex (a1) {\(\overline b\)};
        	    \vertex[right=1cm of a1] (a2);
        	    \vertex[right=1cm of a2] (a3);
        	    \vertex[right=1cm of a3] (a4) {\(b\)};
        	    \vertex[right=1cm of a4] (a5);
        	    \vertex[right=2cm of a5] (a6) {\(u\)};

        	    \vertex[below=2em of a1] (b1) {\(d\)};
        	    \vertex[right=1cm of b1] (b2);
        	    \vertex[right=1cm of b2] (b3);
        	    \vertex[right=1cm of b3] (b4) {\(\overline d\)};
        	    \vertex[below=2em of a6] (b5) {\(\overline d\)};

        	    \vertex[above=of a6] (c1) {\(\overline u\)};
        	    \vertex[above=2em of c1] (c3) {\(d\)};
        	    \vertex at ($(c1)!0.5!(c3) - (1cm, 0)$) (c2);

        	    \diagram* {
        	      {[edges=fermion]
        	        (b1) -- (b2) -- (a2) -- (a1),
        	        (b5) -- (b4) -- (b3) -- (a3) -- (a4) -- (a5) -- (a6),
        	      },
        	      (a2) -- [boson, edge label=\(W\)] (a3),
        	      (b2) -- [boson, edge label'=\(W\)] (b3),

        	      (c1) -- [fermion, out=180, in=-45] (c2) -- [fermion, out=45, in=180] (c3),
        	      (a5) -- [boson, bend left, edge label=\(W^{-}\)] (c2),
        	    };

        	    \draw [decoration={brace}, decorate] (b1.south west) -- (a1.north west)
        	          node [pos=0.5, left] {\(B^{0}\)};
        	    \draw [decoration={brace}, decorate] (c3.north east) -- (c1.south east)
        	          node [pos=0.5, right] {\(\pi^{-}\)};
        	    \draw [decoration={brace}, decorate] (a6.north east) -- (b5.south east)
        	          node [pos=0.5, right] {\(\pi^{+}\)};
            \end{feynman}
        \end{tikzpicture}
    \end{center}
\end{onlyenv}
\end{frame}






\section{Circuitos}

\begin{frame}
\frametitle{Circuitos}

\begin{center}
    \lmr
    \begin{tikzpicture}
        \draw (0,0) to[R=\qty{2}{\ohm}, i=?, v<=\qty{84}{\volt}] (2,0) -- (2,2) to[V<=\qty{84}{\volt}] (0,2) -- (0,0);
    \end{tikzpicture}
\end{center}
\end{frame}





\begin{frame}[fragile]
\frametitle{Circuitos}

\begin{minted}{latex}
    \usepackage{circuitikz}
\end{minted}

\pause

\begin{columns}
    \begin{column}{0.75\textwidth}
        \begin{block}{INFORMACIÓN}
            En la documentación de {\lmr Circui\TikZ} utilizan el entorno \texttt{circuitikz}.

            Este entorno es un alias del entorno \texttt{tikzpicture}.
        \end{block}
    \end{column}
\end{columns}
\end{frame}







\begin{frame}[fragile]
\frametitle{Circuitos}

\begin{onlyenv}<1>
    \begin{minted}[breaklines]{latex}
\begin{tikzpicture}
    \draw (0,0) to[isource] (0,3) -- (2,3) to[R] (2,0) -- (0,0);
\end{tikzpicture}
    \end{minted}
\end{onlyenv}

\begin{onlyenv}<2>
    \lmr
    \begin{center}
        \begin{tikzpicture}
            \begin{feynman}[large, vertical=e to f]
                \draw (0,0) to[isource] (0,3) -- (2,3) to[R] (2,0) -- (0,0);
            \end{feynman}
        \end{tikzpicture}
    \end{center}
\end{onlyenv}
\end{frame}






\begin{frame}[fragile]
\frametitle{Circuitos}

\begin{onlyenv}<1>
    \begin{minted}[breaklines]{latex}
\begin{tikzpicture}[american]
    \draw (0,0) to[isource, l=$I_0$] (0,3) -- (2,3) to[R=$R_1$] (2,0) -- (0,0);
    \draw (2,3) -- (4,3) to[R=$R_2$] (4,0) -- (2,0);
\end{tikzpicture}
    \end{minted}
\end{onlyenv}

\begin{onlyenv}<2>
    \lmr
    \begin{center}
        \begin{tikzpicture}[american]
            \draw (0,0) to[isource, l=$I_0$] (0,3) -- (2,3) to[R=$R_1$] (2,0) -- (0,0);
            \draw (2,3) -- (4,3) to[R=$R_2$] (4,0) -- (2,0);
        \end{tikzpicture}
    \end{center}
\end{onlyenv}
\end{frame}





\begin{frame}[fragile]
\frametitle{Circuitos}

\begin{onlyenv}<1>
    \begin{minted}[breaklines]{latex}
\begin{tikzpicture}[american]
    \draw (0,0)
        to[isource, l=$I_0$] (0,3)
        to[short, -*, i=$I_0$] (2,3)
        to[R=$R_1$, i=$i_1$] (2,0)
        -- (0,0);
    \draw (2,3)
        -- (4,3)
        to[R=$R_2$, i=$i_2$] (4,0)
        to[short, -*] (2,0);
\end{tikzpicture}
    \end{minted}
\end{onlyenv}

\begin{onlyenv}<2>
    \lmr
    \begin{center}
        \begin{tikzpicture}[american]
            \draw (0,0)
                to[isource, l=$I_0$] (0,3)
                to[short, -*, i=$I_0$] (2,3)
                to[R=$R_1$, i=$i_1$] (2,0)
                -- (0,0);
            \draw (2,3)
                -- (4,3)
                to[R=$R_2$, i=$i_2$] (4,0)
                to[short, -*] (2,0);
        \end{tikzpicture}
    \end{center}
\end{onlyenv}
\end{frame}






\begin{frame}[fragile]
\frametitle{Circuitos}

\begin{onlyenv}<1>
    \begin{minted}[breaklines]{latex}
\begin{tikzpicture}[american]
    \draw (0,0)
        to[isource, l=$I_0$] (0,3)
        to[short, -*, i=$I_0$] (2,3)
        to[R=$R_1$, i>_=$i_1$] (2,0)
        -- (0,0);
    \draw (2,3)
        -- (4,3)
        to[R=$R_2$, i>_=$i_2$] (4,0)
        to[short, -*] (2,0);
\end{tikzpicture}
    \end{minted}
\end{onlyenv}

\begin{onlyenv}<2>
    \lmr
    \begin{center}
        \begin{tikzpicture}[american]
            \draw (0,0)
                to[isource, l=$I_0$] (0,3)
                to[short, -*, i=$I_0$] (2,3)
                to[R=$R_1$, i>_=$i_1$] (2,0)
                -- (0,0);
            \draw (2,3)
                -- (4,3)
                to[R=$R_2$, i>_=$i_2$] (4,0)
                to[short, -*] (2,0);
        \end{tikzpicture}
    \end{center}
\end{onlyenv}
\end{frame}






\begin{frame}[fragile]
\frametitle{Circuitos}

\begin{onlyenv}<1>
    \begin{minted}[breaklines]{latex}
\begin{tikzpicture}[american, voltage shift=0.5]
    \draw (0,0)
        to[isource, l=$I_0$, v=$V_0$] (0,3)
        to[short, -*, i=$I_0$] (2,3)
        to[R=$R_1$, i>_=$i_1$] (2,0)
        -- (0,0);
    \draw (2,3)
        -- (4,3)
        to[R=$R_2$, i>_=$i_2$] (4,0)
        to[short, -*] (2,0);
\end{tikzpicture}
    \end{minted}
\end{onlyenv}

\begin{onlyenv}<2>
    \lmr
    \begin{center}
        \begin{tikzpicture}[american, voltage shift=0.5]
            \draw (0,0)
                to[isource, l=$I_0$, v=$V_0$] (0,3)
                to[short, -*, i=$I_0$] (2,3)
                to[R=$R_1$, i>_=$i_1$] (2,0)
                -- (0,0);
            \draw (2,3) -- (4,3)
                to[R=$R_2$, i>_=$i_2$] (4,0)
                to[short, -*] (2,0);
        \end{tikzpicture}
    \end{center}
\end{onlyenv}
\end{frame}






\begin{frame}[fragile]
\frametitle{Circuitos}

\begin{onlyenv}<1>
    \begin{minted}[breaklines]{latex}
\begin{tikzpicture}[european, voltage shift=0.5]
    \draw (0,0)
        to[isourceC, l=$I_0$, v=$V_0$] (0,3)
        to[short, -*, f=$I_0$] (2,3)
        to[R=$R_1$, f>_=$i_1$] (2,0)
        -- (0,0);
    \draw (2,3)
        -- (4,3)
        to[R=$R_2$, f>_=$i_2$] (4,0)
        to[short, -*] (2,0);
\end{tikzpicture}
    \end{minted}
\end{onlyenv}

\begin{onlyenv}<2>
    \lmr
    \begin{center}
        \begin{tikzpicture}[european, voltage shift=0.5]
            \draw (0,0)
                to[isourceC, l=$I_0$, v=$V_0$] (0,3)
                to[short, -*, f=$I_0$] (2,3)
                to[R=$R_1$, f>_=$i_1$] (2,0)
                -- (0,0);
            \draw (2,3)
                -- (4,3)
                to[R=$R_2$, f>_=$i_2$] (4,0)
                to[short, -*] (2,0);
        \end{tikzpicture}
    \end{center}
\end{onlyenv}
\end{frame}






\begin{frame}[fragile]
\frametitle{Circuitos}

\begin{onlyenv}<1>
    \begin{minted}[breaklines]{latex}
\begin{tikzpicture}[american, voltage shift=0.5]
    \draw (0,0)
        to[isource, l=$I_0$, v=$V_0$] (0,3)
        to[short, -*, f=$I_0$] (2,3)
        to[R=$R_1$, f>_=$i_1$] (2,0)
        -- (0,0);
    \draw (2,3)
        -- (4,3)
        to[R=$R_2$, f>_=$i_2$] (4,0)
        to[short, -*] (2,0);
    \draw[red, thick] (0.5,2.0) rectangle (4.3,3.9)
        node[pos=0.5, above] {KCL};
\end{tikzpicture}
    \end{minted}
\end{onlyenv}

\begin{onlyenv}<2>
    \lmr
    \begin{center}
        \begin{tikzpicture}[american, voltage shift=0.5]
            \draw (0,0)
                to[isource, l=$I_0$, v=$V_0$] (0,3)
                to[short, -*, f=$I_0$] (2,3)
                to[R=$R_1$, f>_=$i_1$] (2,0)
                -- (0,0);
            \draw (2,3)
                -- (4,3)
                to[R=$R_2$, f>_=$i_2$] (4,0)
                to[short, -*] (2,0);
            \draw[red, thick] (0.5,2.0) rectangle (4.3,3.9)
                node[pos=0.5, above] {KCL};
        \end{tikzpicture}
    \end{center}
\end{onlyenv}
\end{frame}





\begin{frame}[fragile]
\frametitle{Circuitos}

\begin{onlyenv}<1>
    \begin{minted}[breaklines]{latex}
\begin{tikzpicture}
    \node [op amp] {\texttt{OA1}};
\end{tikzpicture}
    \end{minted}
\end{onlyenv}

\begin{onlyenv}<2>
    \lmr
    \begin{center}
        \begin{tikzpicture}
            \node [op amp] {\texttt{OA1}};
        \end{tikzpicture}
    \end{center}
\end{onlyenv}
\end{frame}





\begin{frame}[fragile]
\frametitle{Circuitos}

\begin{onlyenv}<1>
    \begin{minted}[breaklines]{latex}
\begin{tikzpicture}
    \node [op amp, noinv input up] {\texttt{OA1}};
\end{tikzpicture}
    \end{minted}
\end{onlyenv}

\begin{onlyenv}<2>
    \lmr
    \begin{center}
        \begin{tikzpicture}
            \node [op amp, noinv input up] {\texttt{OA1}};
        \end{tikzpicture}
    \end{center}
\end{onlyenv}
\end{frame}





\begin{frame}[fragile]
\frametitle{Circuitos}

\begin{center}
    \begin{tikzpicture}
        \node [op amp] (OA) {\texttt{OA1}};

        \pause

        \begin{scope}[red,fill=transparent]
            \draw (OA.-) node[circle,red,draw,inner sep=1pt,pin={[red,overlay,pin edge={red, overlay}]left:$-$}] {};
            \draw (OA.+) node[circle,red,draw,inner sep=1pt,pin={[red,overlay,pin edge={red, overlay}]left:$+$}] {};
            \draw (OA.up) node[circle,red,draw,inner sep=1pt,pin={[red,overlay,pin edge={red, overlay}]above:up}] {};
            \draw (OA.down) node[circle,red,draw,inner sep=1pt,pin={[red,overlay,pin edge={red, overlay}]below:down}] {};
            \draw (OA.out) node[circle,red,draw,inner sep=1pt,pin={[red,overlay,pin edge={red, overlay}]right:out}] {};
        \end{scope}
    \end{tikzpicture}
\end{center}
\end{frame}





\begin{frame}[fragile]
\frametitle{Circuitos}

\begin{onlyenv}<1>
    \begin{minted}[breaklines]{latex}
\begin{tikzpicture}
    \node [op amp, noinv input up] (OA) {\texttt{OA1}};

    \draw (OA.+) to[short, -o] ++(-1,0) node [above] {$v_i$};
    \draw (OA.-)
        -- ++(0,-1)
        coordinate (FB)
        to[R=$R_1$] ++(0,-2)
        node [ground] {};
\end{tikzpicture}
    \end{minted}
\end{onlyenv}

\begin{onlyenv}<2>
    \lmr
    \begin{center}
        \begin{tikzpicture}
            \node [op amp, noinv input up] (OA) {\texttt{OA1}};

            \draw (OA.+) to[short, -o] ++(-1,0) node [above] {$v_i$};
            \draw (OA.-)
                -- ++(0,-1)
                coordinate (FB)
                to[R=$R_1$] ++(0,-2)
                node [ground] {};
        \end{tikzpicture}
    \end{center}
\end{onlyenv}
\end{frame}





\begin{frame}[fragile]
\frametitle{Circuitos}

\begin{onlyenv}<1>
    \begin{minted}[breaklines]{latex}
\begin{tikzpicture}
    \node [op amp, noinv input up] (OA) {\texttt{OA1}};

    \draw (OA.+) to[short, -o] ++(-1,0) node [above] {$v_i$};
    \draw (OA.-)
        -- ++(0,-1)
        coordinate (FB)
        to[R=$R_1$] ++(0,-2)
        node [ground] {};
    \draw (FB)
        to[R=$R_2$,*-] (FB -| OA.out)
        -- (OA.out)
        to[short,*-o] ++(1,0)
        node [above] {$v_o$};
\end{tikzpicture}
    \end{minted}
\end{onlyenv}

\begin{onlyenv}<2>
    \lmr
    \begin{center}
        \begin{tikzpicture}
            \node [op amp, noinv input up] (OA) {\texttt{OA1}};

            \draw (OA.+) to[short, -o] ++(-1,0) node [above] {$v_i$};
            \draw (OA.-)
                -- ++(0,-1)
                coordinate (FB)
                to[R=$R_1$] ++(0,-2)
                node [ground] {};
            \draw (FB)
                to[R=$R_2$,*-] (FB -| OA.out)
                -- (OA.out)
                to[short,*-o] ++(1,0)
                node [above] {$v_o$};
        \end{tikzpicture}
    \end{center}
\end{onlyenv}
\end{frame}





\begin{frame}[fragile]
\frametitle{Circuitos}

\begin{onlyenv}<1>
    \begin{minted}[breaklines]{latex}
\ctikzset{
    amplifiers/fill=red!80,
    component text=left,
}
    \end{minted}
\end{onlyenv}

\begin{onlyenv}<2>
    \lmr
    \begin{center}
        \ctikzset{amplifiers/fill=red!80, component text=left}
        \begin{tikzpicture}
            \node [op amp, noinv input up] (OA) {\texttt{OA1}};

            \draw (OA.+) to[short, -o] ++(-1,0) node [above] {$v_i$};
            \draw (OA.-)
                -- ++(0,-1)
                coordinate (FB)
                to[R=$R_1$] ++(0,-2)
                node [ground] {};
            \draw (FB)
                to[R=$R_2$,*-] (FB -| OA.out)
                -- (OA.out)
                to[short,*-o] ++(1,0)
                node [above] {$v_o$};
        \end{tikzpicture}
    \end{center}
\end{onlyenv}
\end{frame}





\begin{frame}[fragile]
\frametitle{Circuitos}

\begin{onlyenv}<1>
    \scriptsize
    \begin{minted}[breaklines]{latex}
\begin{tikzpicture}[american]
    \node [nmos] (Q1) {Q1};
    \draw (Q1.S)
        to[R, l2^=$R_S$ and \qty{5}{\kilo\ohm}] ++(0,-3)
        node [vee] (VEE) {$V_{EE}=\qty{-10}{\volt}$};
    \draw (Q1.D)
        to[R, l2_=$R_D$ and \qty{10}{\kilo\ohm}] ++(0,3)
        node [vcc] (VCC) {$V_{CC}=\qty{10}{\volt}$};
    \draw (Q1.S)
        to[short] ++(2,0)
        to[C=$C_1$] ++(0,-1.5)
        node [ground] (GND) {};
    \draw (Q1.G)
        to[short] ++(-1,0) coordinate (in)
        to[R, l2^=$R_G$ and \qty{1}{\mega\ohm}] (in |- GND)
        node [ground] {};
    \draw (in)
        to[C, l_=$C_2$,*-o] ++(-1.5,0)
        node [left] (vi1) {$v_i=v_{i1}$};
\end{tikzpicture}
    \end{minted}
\end{onlyenv}

\begin{onlyenv}<2>
    \lmr
    \begin{center}
        \begin{tikzpicture}[american,scale=0.75,transform shape]
            \node [nmos] (Q1) {Q1};
            \draw (Q1.S)
                to[R, l2^=$R_S$ and \qty{5}{\kilo\ohm}] ++(0,-3)
                node [vee] (VEE) {$V_{EE}=\qty{-10}{\volt}$};
            \draw (Q1.D)
                to[R, l2_=$R_D$ and \qty{10}{\kilo\ohm}] ++(0,3)
                node [vcc] (VCC) {$V_{CC}=\qty{10}{\volt}$};
            \draw (Q1.S)
                to[short] ++(2,0)
                to[C=$C_1$] ++(0,-1.5)
                node [ground] (GND) {};
            \draw (Q1.G)
                to[short] ++(-1,0) coordinate (in)
                to[R, l2^=$R_G$ and \qty{1}{\mega\ohm}] (in |- GND)
                node [ground] {};
            \draw (in)
                to[C, l_=$C_2$,*-o] ++(-1.5,0)
                node [left] (vi1) {$v_i=v_{i1}$};
        \end{tikzpicture}
    \end{center}
\end{onlyenv}
\end{frame}





\begin{frame}[fragile]
\frametitle{Circuitos}
\framesubtitle{Ejercicio}

\begin{onlyenv}<1>
    \lmr
    \begin{center}
        \begin{tikzpicture}[american]
            \draw (0,0)
                to[battery,v_=$V$,i=$I$] (0,-3)
                to[capacitor,l_=$C$] (3,-3)
                to[inductor,l_=$L$] (3,0)
                to[resistor,l_=$R$] (0,0);
        \end{tikzpicture}
    \end{center}
\end{onlyenv}

\begin{onlyenv}<2>
    \begin{minted}[breaklines]{latex}
\begin{tikzpicture}[american]
    \draw (0,0)
        to[battery,v_=$V$,i=$I$] (0,-3)
        to[capacitor,l_=$C$] (3,-3)
        to[inductor,l_=$L$] (3,0)
        to[resistor,l_=$R$] (0,0);
\end{tikzpicture}
    \end{minted}
\end{onlyenv}
\end{frame}






\begin{frame}[fragile]
\frametitle{Circuitos}

\begin{columns}
    \begin{column}{0.3\textwidth}
        \centering
        \textbf{Dual-in-Line Package Chip}\bigskip

        \begin{tikzpicture}
            \node [dipchip] {};
        \end{tikzpicture}\par
    \end{column}

    \begin{column}{0.3\textwidth}
        \centering
        \textbf{Quad-Flat Package Chip}\bigskip

        \begin{tikzpicture}
            \node [qfpchip] {};
        \end{tikzpicture}\par
    \end{column}
\end{columns}
\end{frame}





\begin{frame}[fragile]
\frametitle{Circuitos}

\begin{onlyenv}<1>
    \begin{minted}[breaklines]{latex}
\begin{tikzpicture}
    \node [
        dipchip,
        num pins=12,
    ] {IC1};
\end{tikzpicture}
    \end{minted}
\end{onlyenv}

\begin{onlyenv}<2>
    \lmr
    \begin{center}
        \begin{tikzpicture}
            \node [
                dipchip,
                num pins=12,
            ] {IC1};
        \end{tikzpicture}
    \end{center}
\end{onlyenv}
\end{frame}





\begin{frame}
\frametitle{Circuitos}

\begin{columns}
    \begin{column}{0.75\textwidth}
        \begin{alertblock}{IMPORTANTE}
            El número de pines que puede recibir \texttt{num pins} debe ser par y mayor que cero.

            En caso de que sea un número impar, {\lmr Circui\TikZ} lo aproximara al par inmediatamente anterior.
        \end{alertblock}
    \end{column}
\end{columns}
\end{frame}





\begin{frame}[fragile]
\frametitle{Circuitos}

\begin{onlyenv}<1>
    \begin{minted}[breaklines]{latex}
\begin{tikzpicture}
    \node [
        dipchip,
        num pins=8,
    ] (555) {555};

    \draw (555.pin 8) -- ++(0.5,0) to[short,-o] ++(0,0.5) node [above] {\qty{5}{\volt}};
    \draw (555.pin 4) to[short,-o] ++(-0.5,0) node [left] {\qty{5}{\volt}};
    \draw (555.pin 1) -- ++(0,0.5) -- ++(-0.5,0) node [ground] {};

    ...
\end{tikzpicture}
    \end{minted}
\end{onlyenv}

\begin{onlyenv}<2>
    \lmr
    \begin{center}
        \begin{tikzpicture}
            \node [
                dipchip,
                num pins=8,
            ] (555) {555};

            \draw (555.pin 8) -- ++(0.5,0) to[short,-o] ++(0,0.5) node [above] {\qty{5}{\volt}};
            \draw (555.pin 4) to[short,-o] ++(-0.5,0) node [left] {\qty{5}{\volt}};
            \draw (555.pin 1) -- ++(0,0.5) -- ++(-0.5,0) node [ground] {};

            \draw (555.pin 3) -- ++(-4,0)
                to [resistor,l_=$\qty{470}{\ohm}$] ++(0,-1.5)
                to [led,l_=LED] ++(0,-1.5)
                node [ground] {};

            \draw (555.pin 2)
                -- ++(-2,0)
                coordinate (TRG)
                to[short,-o] ++(0,2)
                node [above] {\qty{5}{\volt}};

            \draw (TRG)
                to[short,*-] (555.south -| TRG)
                to[short,-*] ++(0,-1)
                coordinate (TRO)
                to[resistor,l=\qty{47}{\kilo\ohm}] ++(0,-1.5)
                node [ground] {};

            \draw (555.pin 6)
                -- ++(0.5,0)
                |- (TRO);
        \end{tikzpicture}
    \end{center}
\end{onlyenv}
\end{frame}






\section{Asymptote}

\begin{frame}[fragile]
\begin{center}
    \begin{asy}
        size(140,80,IgnoreAspect);

        picture logo(pair s=0, pen q)
        {
            picture pic;
            pen p=linewidth(2)+fontsize(24pt)+q;
            real a=-0.4;
            real b=0.95;
            real y1=-5;
            real y2=-3y1/2;
            path A=(a,0){dir(10)}::{dir(89.5)}(0,y2);
            draw(pic,A,p);
            draw(pic,(0,y1){dir(88.3)}::{dir(20)}(b,0),p);
            real c=0.5*a;
            pair z=(0,2.5);
            label(pic,"{\it symptote}",z,0.25*E+0.169S,p);
            pair w=(0,1.7);
            draw(pic,intersectionpoint(A,w-1--w)--w,p);
            draw(pic,(0,y1)--(0,y2),p);
            draw(pic,(a,0)--(b,0),p);
            return shift(s)*pic;
        }

        pair z=(-0.015,0.08);
        for(int x=0; x < 10; ++x)
            add(logo(0.1*x*z,gray(0.04*x)));

        add(logo(red));
    \end{asy}
\end{center}
\end{frame}





\begin{frame}
\frametitle{¿Qué es Asymptote?}

\begin{itemize}
    \item {\lmr\itshape Asymptote} es un sistema de composición de gráficos vectoriales para dibujo técnico.
    \item Permite el manejo de elementos en tres dimensiones.
    \item Posee una capa de integración con {\lmr\LaTeX}.
\end{itemize}
\end{frame}





\begin{frame}
\frametitle{Prerrequisitos}

\begin{itemize}
    \item Compilador de {\lmr\itshape Asymptote}
    \item Distribución de {\lmr\LaTeX}
    \item GhostScript
    \item ImageMagick
    \item Visor de PostScript
\end{itemize}
\end{frame}






\begin{frame}[fragile]
\frametitle{Asymptote}

\begin{onlyenv}<1>
    \begin{minted}[breaklines]{latex}
\usepackage{asymptote}

\begin{asy}
    settings.prc = false;
    import three;

    size(4cm,0);

    draw((0,0,0)--(2,0,0), blue);
    draw((0,0,0)--(0,2,0), green);
    draw((0,0,0)--(0,0,2), red);
\end{asy}
    \end{minted}
\end{onlyenv}

\begin{onlyenv}<2>
    \lmr
    \begin{center}
        \begin{asy}
            settings.prc = false;
            import three;
            currentlight.background=opacity(0.0);
            currentpen = white;

            size(4cm,0);

            draw((0,0,0)--(2,0,0), blue);
            draw((0,0,0)--(0,2,0), green);
            draw((0,0,0)--(0,0,2), red);
        \end{asy}
    \end{center}
\end{onlyenv}
\end{frame}







\begin{frame}[fragile]
\frametitle{Asymptote}

\begin{onlyenv}<1>
    \begin{minted}[breaklines]{asymptote}
settings.prc = false;
import three;

size(4cm,0);

draw((0,0,0)--(2,0,0), blue);
draw((0,0,0)--(0,2,0), green);
draw((0,0,0)--(0,0,2), red);
    \end{minted}
\end{onlyenv}

\begin{onlyenv}<2>
    \lmr
    \begin{center}
        \begin{asy}
            settings.prc = false;
            import three;
            currentlight.background=opacity(0.0);
            currentpen = white;

            size(4cm,0);

            draw((0,0,0)--(2,0,0), blue);
            draw((0,0,0)--(0,2,0), green);
            draw((0,0,0)--(0,0,2), red);
        \end{asy}
    \end{center}
\end{onlyenv}
\end{frame}







\begin{frame}[fragile]
\frametitle{Asymptote}

\begin{onlyenv}<1>
    \begin{minted}[breaklines]{asymptote}
settings.prc=false;
import three;

size(4cm,0);

draw((0,0,0)--(2,0,0), blue);
draw((0,0,0)--(0,2,0), green);
draw((0,0,0)--(0,0,2), red);
draw(box(O, (0.5, 1.5, 1)), blue);
    \end{minted}
\end{onlyenv}

\begin{onlyenv}<2>
    \lmr
    \begin{center}
        \begin{asy}
            settings.prc=false;
            import three;
            currentlight.background=opacity(0.0);
            currentpen = white;

            size(4cm,0);

            draw((0,0,0)--(2,0,0), blue);
            draw((0,0,0)--(0,2,0), green);
            draw((0,0,0)--(0,0,2), red);
            draw(box(O, (0.5, 1.5, 1)), blue);
        \end{asy}
    \end{center}
\end{onlyenv}
\end{frame}






\begin{frame}[fragile]
\frametitle{Asymptote}

\begin{onlyenv}<1>
    \begin{minted}[breaklines]{asymptote}
settings.prc=false;
settings.render=16;
import three;

size(3.55cm, 0);

draw(-1.5X -- 1.5X, arrow=Arrow3(TeXHead2), L=Label("$x$", position=EndPoint, align=W));
draw(-1.5Y -- 1.5Y, arrow=Arrow3(TeXHead2), L=Label("$y$", position=EndPoint));
draw(-1.5Z -- 1.5Z, arrow=Arrow3(TeXHead2), L=Label("$z$", position=EndPoint));
draw(unitsphere, surfacepen = material(white, emissivepen = gray(0.2)));
    \end{minted}
\end{onlyenv}

\begin{onlyenv}<2>
    \lmr
    \begin{center}
        \begin{asy}
            settings.prc=false;
            settings.render=16;
            import three;
            currentlight.background=opacity(0.0);
            currentpen = white;

            size(3.55cm, 0);

            draw(-1.5X -- 1.5X, arrow=Arrow3(TeXHead2), L=Label("$x$", position=EndPoint, align=W));
            draw(-1.5Y -- 1.5Y, arrow=Arrow3(TeXHead2), L=Label("$y$", position=EndPoint));
            draw(-1.5Z -- 1.5Z, arrow=Arrow3(TeXHead2), L=Label("$z$", position=EndPoint));
            draw(unitsphere, surfacepen = material(white, emissivepen = gray(0.2)));
        \end{asy}
    \end{center}
\end{onlyenv}
\end{frame}





\begin{frame}[fragile]
\frametitle{Asymptote}

\begin{onlyenv}<1>
    \tiny
    \begin{minted}[breaklines]{asymptote}
settings.render = 8;
defaultpen(fontsize(10pt));
import graph;
import three;

size(4.15cm, 0);

currentprojection = orthographic(5,0,10, up=Y);
real xmin = -0.1;
real xmax = 2;
real ymin = -0.1;
real ymax = 2;
real margin = 0.2;
real f(real x) { return sqrt(x); }
path s = graph(f, 0, 2, operator..);
path3 p3 = path3(s);
draw(p3);
surface solidsurface = surface(p3, c=O, axis=X);
draw(solidsurface, white);
draw(xmin*X -- xmax*X);
draw(ymin*Y -- ymax*Y);
path planeoutline = box((xmin, ymin), (xmax+margin, ymax+margin));
draw(surface(planeoutline), surfacepen=white);
    \end{minted}
\end{onlyenv}

\begin{onlyenv}<2>
    \lmr
    \begin{center}
        \begin{asy}
            settings.render = 8;
            defaultpen(fontsize(10pt));
            import graph;
            import three;
            currentlight.background=opacity(0.0);
            currentpen = white;

            size(4.15cm, 0);

            currentprojection = orthographic(5,0,10, up=Y);
            real xmin = -0.1;
            real xmax = 2;
            real ymin = -0.1;
            real ymax = 2;
            real margin = 0.2;
            real f(real x) { return sqrt(x); }
            path s = graph(f, 0, 2, operator..);
            path3 p3 = path3(s);
            draw(p3);
            surface solidsurface = surface(p3, c=O, axis=X);
            draw(solidsurface, white);
            draw(xmin*X -- xmax*X);
            draw(ymin*Y -- ymax*Y);
            path planeoutline = box((xmin, ymin), (xmax+margin, ymax+margin));
            draw(surface(planeoutline), surfacepen=white);
        \end{asy}
    \end{center}
\end{onlyenv}
\end{frame}







\begin{frame}[allowframebreaks]
\frametitle{Referencias}
\nocite{*}
\printbibliography
\end{frame}
\end{document}
